
\documentclass{article}
\usepackage{colortbl}
\usepackage{makecell}
\usepackage{multirow}
\usepackage{supertabular}

\begin{document}

\newcounter{utterance}

\centering \large Interaction Transcript for game `function\_detective', experiment `string\_passive\_labeled\_pairs', episode 159 with Qwen3.6{-}35B{-}A3B{-}FP8{-}without{-}reasoning.
\vspace{24pt}

{ \footnotesize  \setcounter{utterance}{1}
\setlength{\tabcolsep}{0pt}
\begin{supertabular}{c@{$\;$}|p{.15\linewidth}@{}p{.15\linewidth}p{.15\linewidth}p{.15\linewidth}p{.15\linewidth}p{.15\linewidth}}
    \# & \multicolumn{2}{c}{Player} && \multicolumn{2}{c}{Game Master} \\
    \hline

    \theutterance \stepcounter{utterance}  
    & & & \multicolumn{4}{p{0.6\linewidth}}{
        \cellcolor[rgb]{0.9,0.9,0.9}{
            \makecell[{{p{\linewidth}}}]{
                \texttt{\tiny{[P1$\langle$GM]}}
                \texttt{I have a hidden Python function for you to figure out.} \\
\\ 
\texttt{Your job is to infer what it does and then submit a correct Python implementation.} \\
\\ 
\texttt{Function signature:} \\
\texttt{(text: str) {-}> str} \\
\\ 
\texttt{You have 15 turns total.} \\
\\ 
\texttt{I reveal candidate input{-}output pairs one at a time and tell you whether each one is correct.} \\
\\ 
\texttt{You must make your best SOLVE attempt no later than the second{-}to{-}last turn.} \\
\texttt{Do not wait until the very end.} \\
\\ 
\texttt{You can use these actions:} \\
\texttt{{-} JSON format described at the end of the prompt} \\
\texttt{{-} SOLVE: ```python ... ```} \\
\\ 
\texttt{Do not add import statements in your guess.} \\
\texttt{Anything you might need is already available, so just write the function body in hypothesis and return the answer.} \\
\\ 
\texttt{Preloaded examples:} \\
\texttt{1. (('   ', '20 20 20'), True)} \\
\\ 
\texttt{Scoring:} \\
\texttt{{-} correctness is required for any non{-}zero score} \\
\texttt{{-} in interactive modes, solving with fewer information requests is better} \\
\\ 
\texttt{Input formatting reminders:} \\
\texttt{{-} strings must be quoted} \\
\texttt{{-} booleans must be True/False} \\
\texttt{{-} use commas between multiple values} \\
\texttt{{-} do not write any extra action tags in the same message} \\
\\ 
\texttt{Format example:} \\
\texttt{You write:} \\
\texttt{\{'current\_mode': <probing/confirming/don't know; choose whichever best describes your current state in the task>, 'input\_rationale': <brief description of why you chose this hypothesis>, 'hypothesis': <either the exact string "unknown" if you have no guess yet OR your current best guess as runnable Python matching the given signature, with a one{-}line docstring stating what it does. Never write placeholder code, if you cannot commit to a guess, use "unknown". Match the given signature.} \\
\texttt{e.g.} \\
\texttt{def f(x):} \\
\texttt{	"""Returns the absolute value of x."""} \\
\texttt{	return x if x >= 0 else {-}x} \\
\\ 
\texttt{Use "unknown" if you have no guess.>, 'confidence\_score': <score between 1{-}100. 100 means highest confidence in your hypothesis>\}} \\
\\ 
\texttt{Output only one valid json per turn.} \\
            }
        }
    }
     \\ \\

    \theutterance \stepcounter{utterance}  
    & \multicolumn{4}{p{0.6\linewidth}}{
        \cellcolor[rgb]{0.561,0.737,0.561}{
            \makecell[{{p{\linewidth}}}]{
                \texttt{\tiny{[P1$\rangle$GM]}}
                \texttt{\{'current\_mode': 'probing', 'input\_rationale': 'I have one preloaded example: input ("   ", "20 20 20") returns True. This suggests the function might be checking if two strings are equal, or if they represent the same value after some normalization. The first string is whitespace, the second is numbers. They are not equal as strings. So it's not a simple equality check. Let me probe with more examples to understand the pattern.', 'hypothesis': 'unknown', 'confidence\_score': 10\}} \\
            }
        }
    }
    & & \\ \\

    \theutterance \stepcounter{utterance}  
    & & & \multicolumn{4}{p{0.6\linewidth}}{
        \cellcolor[rgb]{0.9,0.9,0.9}{
            \makecell[{{p{\linewidth}}}]{
                \texttt{\tiny{[P1$\langle$GM]}}
                \texttt{Game Violation: ParseError: Response is not valid JSON.} \\
\texttt{Output ONLY:} \\
\texttt{One JSON with confidence\_score, current\_mode, input\_rationale, and hypothesis keys ...} \\
\texttt{OR} \\
\texttt{SOLVE: ```python ... ```} \\
            }
        }
    }
     \\ \\

    \theutterance \stepcounter{utterance}  
    & \multicolumn{4}{p{0.6\linewidth}}{
        \cellcolor[rgb]{0.561,0.737,0.561}{
            \makecell[{{p{\linewidth}}}]{
                \texttt{\tiny{[P1$\rangle$GM]}}
                \texttt{\{"current\_mode": "probing", "input\_rationale": "I have one preloaded example: input ('   ', '20 20 20') returns True. This is strange because the strings are clearly different. Perhaps the function checks if the strings are equal after stripping whitespace? '   '.strip() is '' and '20 20 20'.strip() is '20 20 20', which are not equal. Maybe it checks if they are equal after some other transformation. Let me try a simple equality check first to see if the preloaded example is misleading or if there's a different interpretation. Actually, let me just probe with a simple case to understand the function better.", "hypothesis": "unknown", "confidence\_score": 10\}} \\
            }
        }
    }
    & & \\ \\

    \theutterance \stepcounter{utterance}  
    & & & \multicolumn{4}{p{0.6\linewidth}}{
        \cellcolor[rgb]{0.9,0.9,0.9}{
            \makecell[{{p{\linewidth}}}]{
                \texttt{\tiny{[P1$\langle$GM]}}
                \texttt{OUTPUT: (('ab', '61 62\_alt'), False)} \\
            }
        }
    }
     \\ \\

    \theutterance \stepcounter{utterance}  
    & \multicolumn{4}{p{0.6\linewidth}}{
        \cellcolor[rgb]{0.561,0.737,0.561}{
            \makecell[{{p{\linewidth}}}]{
                \texttt{\tiny{[P1$\rangle$GM]}}
                \texttt{\{"current\_mode": "probing", "input\_rationale": "The preloaded example ('   ', '20 20 20') returned True. The new example ('ab', '61 62\_alt') returned False. Let's analyze the new example. 'ab' has ASCII values 97 and 98. '61 62' are hex values for 'a' and 'b'. So 'ab' corresponds to hex '61 62'. The second string is '61 62\_alt'. The '\_alt' part might be a suffix or modifier. Since the result is False, maybe the function checks if the first string matches the hex{-}decoded content of the second string, ignoring any suffix? But the preloaded example ('   ', '20 20 20') returned True. '   ' is three spaces. ASCII 32 is space. Hex for 32 is '20'. So '   ' corresponds to '20 20 20'. This matches perfectly. So the function seems to check if the first string, when converted to hex (space{-}separated), equals the second string (possibly ignoring a suffix like '\_alt'). Let's test this hypothesis. For 'ab', hex is '61 62'. The second string is '61 62\_alt'. If we strip '\_alt', we get '61 62',} \\
            }
        }
    }
    & & \\ \\

    \theutterance \stepcounter{utterance}  
    & & & \multicolumn{4}{p{0.6\linewidth}}{
        \cellcolor[rgb]{0.9,0.9,0.9}{
            \makecell[{{p{\linewidth}}}]{
                \texttt{\tiny{[P1$\langle$GM]}}
                \texttt{Game Violation: ParseError: Response is not valid JSON.} \\
\texttt{Output ONLY:} \\
\texttt{One JSON with confidence\_score, current\_mode, input\_rationale, and hypothesis keys ...} \\
\texttt{OR} \\
\texttt{SOLVE: ```python ... ```} \\
            }
        }
    }
     \\ \\

    \theutterance \stepcounter{utterance}  
    & \multicolumn{4}{p{0.6\linewidth}}{
        \cellcolor[rgb]{0.561,0.737,0.561}{
            \makecell[{{p{\linewidth}}}]{
                \texttt{\tiny{[P1$\rangle$GM]}}
                \texttt{\{"current\_mode": "probing", "input\_rationale": "Based on the preloaded example ('   ', '20 20 20') {-}> True and the new example ('ab', '61 62\_alt') {-}> False, I hypothesize that the function checks if the first string, when converted to its hexadecimal representation (space{-}separated), matches the second string after removing any suffix like '\_alt'. For '   ', hex is '20 20 20', which matches '20 20 20'. For 'ab', hex is '61 62', which does not match '61 62\_alt' even if we remove '\_alt' because the comparison might be strict or the suffix handling is different. Let me test with a case where the hex matches exactly without suffix to confirm.", "hypothesis": "unknown", "confidence\_score": 30\}} \\
            }
        }
    }
    & & \\ \\

    \theutterance \stepcounter{utterance}  
    & & & \multicolumn{4}{p{0.6\linewidth}}{
        \cellcolor[rgb]{0.9,0.9,0.9}{
            \makecell[{{p{\linewidth}}}]{
                \texttt{\tiny{[P1$\langle$GM]}}
                \texttt{OUTPUT: (('a{-}b', '61 2d 62'), True)} \\
            }
        }
    }
     \\ \\

    \theutterance \stepcounter{utterance}  
    & \multicolumn{4}{p{0.6\linewidth}}{
        \cellcolor[rgb]{0.561,0.737,0.561}{
            \makecell[{{p{\linewidth}}}]{
                \texttt{\tiny{[P1$\rangle$GM]}}
                \texttt{\{"current\_mode": "probing", "input\_rationale": "The example ('a{-}b', '61 2d 62') returns True. 'a' is 0x61, '{-}' is 0x2d, 'b' is 0x62. So the hex representation of 'a{-}b' is '61 2d 62', which matches the second string exactly. This confirms that the function checks if the space{-}separated hex representation of the first string equals the second string. The previous example ('ab', '61 62\_alt') returned False because '61 62' != '61 62\_alt'. The preloaded example ('   ', '20 20 20') returned True because '20 20 20' == '20 20 20'. So the hypothesis is: the function returns True if the hex representation of the first string (with spaces between each byte) equals the second string, and False otherwise.", "hypothesis": "def f(text: str) {-}> str:\textbackslash{}n    \textbackslash{}"\textbackslash{}"\textbackslash{}"Returns 'True' if the space{-}separated hex representation of the first string equals the second string, else 'False'.\textbackslash{}"\textbackslash{}"\textbackslash{}"\textbackslash{}n    \# The function signature says (text: str) {-}> str, but the examples show two inputs.\textbackslash{}n    \# Wait, the signature is (text: str) {-}> str.} \\
            }
        }
    }
    & & \\ \\

    \theutterance \stepcounter{utterance}  
    & & & \multicolumn{4}{p{0.6\linewidth}}{
        \cellcolor[rgb]{0.9,0.9,0.9}{
            \makecell[{{p{\linewidth}}}]{
                \texttt{\tiny{[P1$\langle$GM]}}
                \texttt{Game Violation: ParseError: Response is not valid JSON.} \\
\texttt{Output ONLY:} \\
\texttt{One JSON with confidence\_score, current\_mode, input\_rationale, and hypothesis keys ...} \\
\texttt{OR} \\
\texttt{SOLVE: ```python ... ```} \\
            }
        }
    }
     \\ \\

    \theutterance \stepcounter{utterance}  
    & \multicolumn{4}{p{0.6\linewidth}}{
        \cellcolor[rgb]{0.561,0.737,0.561}{
            \makecell[{{p{\linewidth}}}]{
                \texttt{\tiny{[P1$\rangle$GM]}}
                \texttt{\{"current\_mode": "probing", "input\_rationale": "I need to clarify the function signature. The prompt says (text: str) {-}> str, but all examples provided have two arguments. This suggests the 'text' parameter might be a tuple or the signature description is misleading. However, looking at the examples: ('   ', '20 20 20') {-}> True, ('ab', '61 62\_alt') {-}> False, ('a{-}b', '61 2d 62') {-}> True. It seems the function takes two strings. Let me assume the signature in the prompt is a simplification or error, and the actual function takes two arguments. Or perhaps 'text' is a single string containing both parts? No, the examples clearly show two separate strings. Let me try to probe with a case where the hex matches exactly to confirm the logic, and also check if the return type is boolean or string. The examples show True/False (booleans). But the signature says {-}> str. This is contradictory. Let me assume the return value is a string 'True' or 'False' or the boolean itself. Given the output format shows True/False, it's likely boolean. But the signature says str. I will assume the function returns a string 'True' or 'False' or the boolean. Let's stick to the logic: check if hex representation matches. I'll probe with a negative case where hex} \\
            }
        }
    }
    & & \\ \\

    \theutterance \stepcounter{utterance}  
    & & & \multicolumn{4}{p{0.6\linewidth}}{
        \cellcolor[rgb]{0.9,0.9,0.9}{
            \makecell[{{p{\linewidth}}}]{
                \texttt{\tiny{[P1$\langle$GM]}}
                \texttt{Game Violation: ParseError: Response is not valid JSON.} \\
\texttt{Output ONLY:} \\
\texttt{One JSON with confidence\_score, current\_mode, input\_rationale, and hypothesis keys ...} \\
\texttt{OR} \\
\texttt{SOLVE: ```python ... ```} \\
            }
        }
    }
     \\ \\

    \theutterance \stepcounter{utterance}  
    & \multicolumn{4}{p{0.6\linewidth}}{
        \cellcolor[rgb]{0.561,0.737,0.561}{
            \makecell[{{p{\linewidth}}}]{
                \texttt{\tiny{[P1$\rangle$GM]}}
                \texttt{\{"current\_mode": "probing", "input\_rationale": "I have strong evidence that the function compares the first string's hex representation (space{-}separated bytes) with the second string. Examples: '   ' {-}> '20 20 20' matches '20 20 20' (True). 'a{-}b' {-}> '61 2d 62' matches '61 2d 62' (True). 'ab' {-}> '61 62' does not match '61 62\_alt' (False). The signature says (text: str) {-}> str, but inputs are pairs. This is confusing. Perhaps the input is a single string formatted as 'str1 str2'? Or maybe the signature is wrong. Let me try to submit a solution assuming the function takes two arguments, as that fits the data. If the signature is strictly one string, I might need to parse it. But 'text: str' usually means one string. Let's look at the preloaded example again: (('   ', '20 20 20'), True). The input is a tuple. The function signature might be misleading or I am misinterpreting. Let's assume the function actually takes two strings based on the examples. I will write a function that takes two strings and returns a boolean. If the system expects a single string input, I might fail. But given the examples,} \\
            }
        }
    }
    & & \\ \\

    \theutterance \stepcounter{utterance}  
    & & & \multicolumn{4}{p{0.6\linewidth}}{
        \cellcolor[rgb]{0.9,0.9,0.9}{
            \makecell[{{p{\linewidth}}}]{
                \texttt{\tiny{[P1$\langle$GM]}}
                \texttt{Game Violation: ParseError: Response is not valid JSON.} \\
\texttt{Output ONLY:} \\
\texttt{One JSON with confidence\_score, current\_mode, input\_rationale, and hypothesis keys ...} \\
\texttt{OR} \\
\texttt{SOLVE: ```python ... ```} \\
            }
        }
    }
     \\ \\

    \theutterance \stepcounter{utterance}  
    & \multicolumn{4}{p{0.6\linewidth}}{
        \cellcolor[rgb]{0.561,0.737,0.561}{
            \makecell[{{p{\linewidth}}}]{
                \texttt{\tiny{[P1$\rangle$GM]}}
                \texttt{\{"current\_mode": "confirming", "input\_rationale": "The pattern is clear: the function takes two strings. It converts the first string to a space{-}separated hexadecimal representation of its bytes. It then checks if this hex string exactly equals the second string. If they match, it returns True; otherwise, False. The signature (text: str) {-}> str in the prompt seems to be a simplification or error, as all examples involve two arguments and return booleans. I will implement the function to take two arguments as per the examples.", "hypothesis": "def f(s1: str, s2: str) {-}> bool:\textbackslash{}n    \textbackslash{}"\textbackslash{}"\textbackslash{}"Returns True if the space{-}separated hex representation of s1 equals s2, else False.\textbackslash{}"\textbackslash{}"\textbackslash{}"\textbackslash{}n    hex\_repr = ' '.join(format(ord(c), '02x') for c in s1)\textbackslash{}n    return hex\_repr == s2", "confidence\_score": 95\}} \\
            }
        }
    }
    & & \\ \\

    \theutterance \stepcounter{utterance}  
    & & & \multicolumn{4}{p{0.6\linewidth}}{
        \cellcolor[rgb]{0.9,0.9,0.9}{
            \makecell[{{p{\linewidth}}}]{
                \texttt{\tiny{[P1$\langle$GM]}}
                \texttt{OUTPUT: (('aA', '61 41\_alt'), False)} \\
            }
        }
    }
     \\ \\

    \theutterance \stepcounter{utterance}  
    & \multicolumn{4}{p{0.6\linewidth}}{
        \cellcolor[rgb]{0.561,0.737,0.561}{
            \makecell[{{p{\linewidth}}}]{
                \texttt{\tiny{[P1$\rangle$GM]}}
                \texttt{\{"current\_mode": "confirming", "input\_rationale": "The example ('aA', '61 41\_alt') returns False. 'a' is 0x61, 'A' is 0x41. So hex is '61 41'. The second string is '61 41\_alt'. Since it returns False, the function does NOT ignore the '\_alt' suffix. It requires an exact match. This confirms my hypothesis: the function checks if the space{-}separated hex representation of the first string exactly equals the second string. The preloaded example ('   ', '20 20 20') {-}> True because '20 20 20' == '20 20 20'. The example ('a{-}b', '61 2d 62') {-}> True because '61 2d 62' == '61 2d 62'. The example ('ab', '61 62\_alt') {-}> False because '61 62' != '61 62\_alt'. The example ('aA', '61 41\_alt') {-}> False because '61 41' != '61 41\_alt'. The logic is solid. I will now submit the solution. Note: The signature in the prompt is (text: str) {-}> str, but the examples clearly show two arguments} \\
            }
        }
    }
    & & \\ \\

    \theutterance \stepcounter{utterance}  
    & & & \multicolumn{4}{p{0.6\linewidth}}{
        \cellcolor[rgb]{0.9,0.9,0.9}{
            \makecell[{{p{\linewidth}}}]{
                \texttt{\tiny{[P1$\langle$GM]}}
                \texttt{Game Violation: ParseError: Response is not valid JSON.} \\
\texttt{Output ONLY:} \\
\texttt{One JSON with confidence\_score, current\_mode, input\_rationale, and hypothesis keys ...} \\
\texttt{OR} \\
\texttt{SOLVE: ```python ... ```} \\
            }
        }
    }
     \\ \\

    \theutterance \stepcounter{utterance}  
    & \multicolumn{4}{p{0.6\linewidth}}{
        \cellcolor[rgb]{0.561,0.737,0.561}{
            \makecell[{{p{\linewidth}}}]{
                \texttt{\tiny{[P1$\rangle$GM]}}
                \texttt{\{"current\_mode": "confirming", "input\_rationale": "The hypothesis is confirmed: the function takes two strings, converts the first to space{-}separated hex, and checks for exact equality with the second. The signature discrepancy is likely a prompt error; the examples dictate the actual behavior. I am confident in this solution.", "hypothesis": "def f(s1: str, s2: str) {-}> bool:\textbackslash{}n    \textbackslash{}"\textbackslash{}"\textbackslash{}"Returns True if the space{-}separated hex representation of s1 equals s2, else False.\textbackslash{}"\textbackslash{}"\textbackslash{}"\textbackslash{}n    hex\_repr = ' '.join(format(ord(c), '02x') for c in s1)\textbackslash{}n    return hex\_repr == s2", "confidence\_score": 95\}} \\
            }
        }
    }
    & & \\ \\

    \theutterance \stepcounter{utterance}  
    & & & \multicolumn{4}{p{0.6\linewidth}}{
        \cellcolor[rgb]{0.9,0.9,0.9}{
            \makecell[{{p{\linewidth}}}]{
                \texttt{\tiny{[P1$\langle$GM]}}
                \texttt{OUTPUT: (('123', '31 32 33'), True)} \\
            }
        }
    }
     \\ \\

    \theutterance \stepcounter{utterance}  
    & \multicolumn{4}{p{0.6\linewidth}}{
        \cellcolor[rgb]{0.561,0.737,0.561}{
            \makecell[{{p{\linewidth}}}]{
                \texttt{\tiny{[P1$\rangle$GM]}}
                \texttt{\{"current\_mode": "confirming", "input\_rationale": "The example ('123', '31 32 33') returns True. '1' is 0x31, '2' is 0x32, '3' is 0x33. Hex is '31 32 33', which matches exactly. This further confirms the hypothesis. I am now ready to submit the final solution.", "hypothesis": "def f(s1: str, s2: str) {-}> bool:\textbackslash{}n    \textbackslash{}"\textbackslash{}"\textbackslash{}"Returns True if the space{-}separated hex representation of s1 equals s2, else False.\textbackslash{}"\textbackslash{}"\textbackslash{}"\textbackslash{}n    hex\_repr = ' '.join(format(ord(c), '02x') for c in s1)\textbackslash{}n    return hex\_repr == s2", "confidence\_score": 100\}} \\
            }
        }
    }
    & & \\ \\

    \theutterance \stepcounter{utterance}  
    & & & \multicolumn{4}{p{0.6\linewidth}}{
        \cellcolor[rgb]{0.9,0.9,0.9}{
            \makecell[{{p{\linewidth}}}]{
                \texttt{\tiny{[P1$\langle$GM]}}
                \texttt{OUTPUT: (('hello world', '68 65 6c 6c 6f 20 77 6f 72 6c 64\_alt'), False)} \\
            }
        }
    }
     \\ \\

    \theutterance \stepcounter{utterance}  
    & \multicolumn{4}{p{0.6\linewidth}}{
        \cellcolor[rgb]{0.561,0.737,0.561}{
            \makecell[{{p{\linewidth}}}]{
                \texttt{\tiny{[P1$\rangle$GM]}}
                \texttt{\{"current\_mode": "confirming", "input\_rationale": "The example ('hello world', '68 65 6c 6c 6f 20 77 6f 72 6c 64\_alt') returns False. 'hello world' hex is '68 65 6c 6c 6f 20 77 6f 72 6c 64'. The second string has '\_alt' suffix, so it doesn't match exactly. This confirms the exact match requirement. I have enough evidence to submit the solution.", "hypothesis": "def f(s1: str, s2: str) {-}> bool:\textbackslash{}n    \textbackslash{}"\textbackslash{}"\textbackslash{}"Returns True if the space{-}separated hex representation of s1 equals s2, else False.\textbackslash{}"\textbackslash{}"\textbackslash{}"\textbackslash{}n    hex\_repr = ' '.join(format(ord(c), '02x') for c in s1)\textbackslash{}n    return hex\_repr == s2", "confidence\_score": 100\}} \\
            }
        }
    }
    & & \\ \\

    \theutterance \stepcounter{utterance}  
    & & & \multicolumn{4}{p{0.6\linewidth}}{
        \cellcolor[rgb]{0.9,0.9,0.9}{
            \makecell[{{p{\linewidth}}}]{
                \texttt{\tiny{[P1$\langle$GM]}}
                \texttt{OUTPUT: (('', ''), True)} \\
            }
        }
    }
     \\ \\

    \theutterance \stepcounter{utterance}  
    & \multicolumn{4}{p{0.6\linewidth}}{
        \cellcolor[rgb]{0.561,0.737,0.561}{
            \makecell[{{p{\linewidth}}}]{
                \texttt{\tiny{[P1$\rangle$GM]}}
                \texttt{\{"current\_mode": "confirming", "input\_rationale": "The example ('', '') returns True. The hex representation of an empty string is an empty string, which matches the second empty string. This is consistent with the hypothesis. All examples support the logic: convert first string to space{-}separated hex, compare exactly with second string. I am ready to submit.", "hypothesis": "def f(s1: str, s2: str) {-}> bool:\textbackslash{}n    \textbackslash{}"\textbackslash{}"\textbackslash{}"Returns True if the space{-}separated hex representation of s1 equals s2, else False.\textbackslash{}"\textbackslash{}"\textbackslash{}"\textbackslash{}n    hex\_repr = ' '.join(format(ord(c), '02x') for c in s1)\textbackslash{}n    return hex\_repr == s2", "confidence\_score": 100\}} \\
            }
        }
    }
    & & \\ \\

    \theutterance \stepcounter{utterance}  
    & & & \multicolumn{4}{p{0.6\linewidth}}{
        \cellcolor[rgb]{0.9,0.9,0.9}{
            \makecell[{{p{\linewidth}}}]{
                \texttt{\tiny{[P1$\langle$GM]}}
                \texttt{OUTPUT: (('a', '61\_alt'), False)} \\
            }
        }
    }
     \\ \\

    \theutterance \stepcounter{utterance}  
    & \multicolumn{4}{p{0.6\linewidth}}{
        \cellcolor[rgb]{0.561,0.737,0.561}{
            \makecell[{{p{\linewidth}}}]{
                \texttt{\tiny{[P1$\rangle$GM]}}
                \texttt{\{"current\_mode": "confirming", "input\_rationale": "The example ('a', '61\_alt') returns False. 'a' is 0x61, so hex is '61'. The second string is '61\_alt', which does not match '61' exactly. This further confirms the exact match requirement. I have sufficient evidence to submit the final solution.", "hypothesis": "def f(s1: str, s2: str) {-}> bool:\textbackslash{}n    \textbackslash{}"\textbackslash{}"\textbackslash{}"Returns True if the space{-}separated hex representation of s1 equals s2, else False.\textbackslash{}"\textbackslash{}"\textbackslash{}"\textbackslash{}n    hex\_repr = ' '.join(format(ord(c), '02x') for c in s1)\textbackslash{}n    return hex\_repr == s2", "confidence\_score": 100\}} \\
            }
        }
    }
    & & \\ \\

    \theutterance \stepcounter{utterance}  
    & & & \multicolumn{4}{p{0.6\linewidth}}{
        \cellcolor[rgb]{0.9,0.9,0.9}{
            \makecell[{{p{\linewidth}}}]{
                \texttt{\tiny{[P1$\langle$GM]}}
                \texttt{OUTPUT: (('hello', '68 65 6c 6c 6f'), True)} \\
            }
        }
    }
     \\ \\

    \theutterance \stepcounter{utterance}  
    & \multicolumn{4}{p{0.6\linewidth}}{
        \cellcolor[rgb]{0.561,0.737,0.561}{
            \makecell[{{p{\linewidth}}}]{
                \texttt{\tiny{[P1$\rangle$GM]}}
                \texttt{\{"current\_mode": "confirming", "input\_rationale": "The example ('hello', '68 65 6c 6c 6f') returns True. 'hello' hex is '68 65 6c 6c 6f', which matches exactly. All examples consistently support the hypothesis: the function takes two strings, converts the first to space{-}separated hex, and checks for exact equality with the second. I am confident in this solution.", "hypothesis": "def f(s1: str, s2: str) {-}> bool:\textbackslash{}n    \textbackslash{}"\textbackslash{}"\textbackslash{}"Returns True if the space{-}separated hex representation of s1 equals s2, else False.\textbackslash{}"\textbackslash{}"\textbackslash{}"\textbackslash{}n    hex\_repr = ' '.join(format(ord(c), '02x') for c in s1)\textbackslash{}n    return hex\_repr == s2", "confidence\_score": 100\}} \\
            }
        }
    }
    & & \\ \\

    \theutterance \stepcounter{utterance}  
    & & & \multicolumn{2}{p{0.3\linewidth}}{
        \cellcolor[rgb]{0.9,0.9,0.9}{
            \makecell[{{p{\linewidth}}}]{
                \texttt{\tiny{[GM$|$GM]}}
                \texttt{Game Verdict: LOSS (turn limit reached)} \\
            }
        }
    }
    & & \\ \\

\end{supertabular}
}

\end{document}
